\section{Ламинирование, бюджет слоя и численные кандидаты}
\label{sec:lam}

Продуктовое правило (предложение~\ref{prop:product}) тратит ширину базы на
ортогональный блок. \emph{Ламинирование} тратит её экономнее: слои сдвигаются
внутрь ячейки базы, и диаметр растёт медленнее, чем в произведении. Ценой
служит верхняя граница диаметра слоёной решётки, которую больше нельзя
выписать в замкнутой форме.

\subsection{Две леммы ламинирования}\label{ssec:lam}

Конструкция строит раскраску $\R^n$ из раскраски $\R^{n-1}$. Пусть
$\Lambda'\subset\R^{n-1}$ --- родительская решётка, $\Gamma'=\ker\psi$ ---
правильная подрешётка индекса~$k$, $c\in\R^{n-1}$ --- \emph{сдвиг слоя},
$t>0$ --- \emph{высота слоя}. Положим
\[
\Lambda=\bigl\{(x+ic,\;it)\ :\ x\in\Lambda',\ i\in\Z\bigr\},\qquad
\Gamma=\bigl\{(x+ic,\;it)\ :\ \psi(x)+ia=0,\ m\mid i\bigr\},
\]
где $a\in\Lambda'/\Gamma'$ --- \emph{склейка}, $m$ --- \emph{модуль слоя};
тогда $[\Lambda:\Gamma]=km$. Как группа $\Lambda=\Lambda'\oplus\Z g$ с
$g=(c,t)$, поэтому $(x,i)\mapsto(\psi(x)+ia,\ i\bmod m)$ --- действительно
характер.

Работу делают два неравенства, каждое --- в одну строку.

\begin{lemma}\label{lem:P1}
$\diam\Lambda\le\sqrt{\diam(\Lambda')^2+t^2}$ при любом сдвиге $c$.
\end{lemma}

\begin{proof}
Для точки $(x,z)$ возьмём ближайший слой $i$ (так что $|z-it|\le t/2$) и
ближайшую к $x-ic$ точку решётки $\Lambda'$; тогда
$\dist^2\bigl((x,z),\Lambda\bigr)\le R(\Lambda')^2+t^2/4$, то есть
$R(\Lambda)^2\le R(\Lambda')^2+t^2/4$.
\end{proof}

\begin{lemma}\label{lem:P2}
$V_0(\Lambda)\subseteq V_0(\Lambda')\times\R$, поэтому для
\emph{горизонтального} вектора $v$ (нулевая слоевая координата, то есть
$v\in\Gamma'$) выполнено $D_\Lambda(v)\ge D_{\Lambda'}(v)$.
\end{lemma}

\begin{proof}
Из $\Lambda'\times\{0\}\subset\Lambda$: если $(x,z)\in V_0(\Lambda)$, то для
всякого $e\in\Lambda'$ имеем $\|x\|^2+z^2\le\|x-e\|^2+z^2$, значит,
$x\in V_0(\Lambda')$. Для горизонтального $v$ точка $v/2$ горизонтальна, и
$\dist\bigl(v/2,V_0(\Lambda')\times\R\bigr)=\dist\bigl(v/2,V_0(\Lambda')\bigr)$.
\end{proof}

Смысл пары лемм: горизонтальные расстояния при подъёме \emph{не убывают}, а
ячейка растёт не более чем на $t^2$ по квадрату диаметра. Поэтому запас
$d(\Gamma')>1$ базовой раскраски --- это в точности бюджет на рост ячейки, а
проверять нужно только векторы с ненулевой слоевой координатой, а таких ---
единицы, поскольку $D(v)\ge\|v\|-\diam$ автоматически закрывает всё длиннее
$2\diam$.

\subsection{Теорема о бюджете слоя}
\label{ssec:layerbudget}

Прямая $t\Z$ --- худшее из возможных покрытий, и естественно заменить её
слоевой \emph{решёткой}.

\begin{proposition}[бюджет слоя]\label{prop:layerbudget}
Пусть база $(\Lambda_0,\Gamma_0)$ имеет ширину $d_0$ и диаметр $\diam_0=2R_0$, и
пусть $\Lambda$ получена ламинированием слоевой решёткой $L\subset\R^m$ с
радиусом покрытия $R_L$ и произвольными поправками. Тогда
$\diam V_0(\Lambda)\le2\sqrt{R_0^2+R_L^2}$ и горизонтальные разделения
сохраняются, поэтому горизонтальная часть условия выполнена, как только
\[
R_L\ \le\ \frac{\diam_0}{2}\sqrt{d_0^2-1}.
\]
\end{proposition}

\begin{proof}
Первое утверждение --- лемма~\ref{lem:P1}, применённая к слоевой решётке
вместо прямой: для точки $(x,z)$ берём ближайший узел $l\in L$ (так что
$\|z-l\|\le R_L$) и ближайшую к $x-c_l$ точку базы, откуда
$R(\Lambda)^2\le R_0^2+R_L^2$. Второе --- лемма~\ref{lem:P2} дословно
($\Lambda_0\times\{0\}\subset\Lambda$). Горизонтальное условие
$D_0\ge\diam V_0(\Lambda)$ при $D_0=d_0\diam_0$ равносильно
$d_0^2\diam_0^2\ge4(R_0^2+R_L^2)=\diam_0^2+4R_L^2$.
\end{proof}

Иначе говоря, \emph{избыток ширины раскраски есть в точности бюджет радиуса
покрытия на дополнительные измерения}. Для $E_8/2401$ ($\diam_0=\sqrt6$,
$d_0^2=7/6$) бюджет равен ровно $R_L\le1/2$. Израсходованный до конца
($R_L=1/2$: прямая единичной высоты либо шестиугольник с
$\lambda_1=\sqrt3/2$), он отвечает вырожденному случаю $d=1$; оба
продуктовых результата работы~\cite{Bounds} (следствие~1 и теорема~6) --- это
два способа потратить его с небольшим остатком: прямая высоты $\sqrt7/3$ даёт $\R^9$ с индексом
$2401\cdot4=9604$ и $\ell=\sqrt{63/61}$, шестиугольник с
$\lambda_1=0{,}8230549$ --- $\R^{10}$ с индексом $2401\cdot19=45619$ и
$\ell=\sqrt{217/214}$. Оптимальные масштабы
доставляет само предложение~\ref{prop:product} ($s_i\propto1/d_i^2$), и в
обоих случаях горизонтальное и слоевое условия связывают одновременно.

Конструкцию полезно делать $\Z[\omega]$-эквивариантной: $\omega$ переставляет три
минимальных направления $A_2$, поэтому вместо трёх независимых слоевых условий
остаётся одна орбита, а свободных параметров --- восемь вместо шестнадцати.

\subsection{Кусочный сертификат диаметра}\label{ssec:cert}

Слабое место всех конструкций этого раздела --- верхняя граница диаметра:
лемма~\ref{lem:P1} строга, но неточна при сдвиге, смещённом с орбит глубоких
дыр, а оценка радиуса покрытия по опорным направлениям сходится \emph{снизу}
и доказательной силы не имеет. Прямые подходы к сертификации (ветвление по
направлениям, перечисление ячеек Делоне) экспоненциальны. Слоёная структура
решётки даёт третий путь.

Для $\Lambda=\{(x+ic,\,it)\}$ расстояние распадается по слоям:
$\dist^2\bigl((x,z),\Lambda\bigr)=\min_i\bigl[f(x-ic)+(z-it)^2\bigr]$, где
$f(y)=\dist^2(y,\Lambda')$. Ограничив минимум слоями $\{0,1\}$ (это лишь
увеличивает правую часть), получаем для $x\in V_0(\Lambda')$, $z\in[0,t]$
задачу с двумя квадратиками $A_0(x)=\|x\|^2$ (на ячейке нулевого слоя) и
$A_1(x)=\|x-c-\mu\|^2$ (на куске $V_0\cap(V_0(\mu)+c)$ общего измельчения двух
сдвинутых разбиений Вороного). Ключевых наблюдений три:

\begin{itemize}[leftmargin=1.4em,itemsep=1pt]
\item кусков конечное число, и они перечислимы: непустота требует
$\|c+\mu\|\le2R(\Lambda')$;
\item разность $A_1-A_0$ \emph{линейна} по $x$, поэтому положение точки
пересечения парабол $z^*=(A_1-A_0+t^2)/(2t)$ относительно $[0,t]$ разрезает
кусок двумя гиперплоскостями на три области, и на каждой максимум по $z$ ---
\emph{выпуклая} квадратика: $A_0+\bigl((A_1-A_0+t^2)/(2t)\bigr)^2$ внутри,
$A_0+t^2$ и $A_1+t^2$ по краям;
\item максимум выпуклой функции по политопу достигается в вершине, а вершины
кусков в размерностях $6$--$8$ перечислимы стандартной машинерией.
\end{itemize}

Итого $R(\Lambda)^2$ ограничен сверху конечным максимумом значений явных
квадратик в явных вершинах --- вычисление, а не оценка. Калибровка: для
индекса $12005$ над $E_8/2401$ (сдвиг в глубокую дыру, точное значение
$R^2=1{,}619025$ известно из точности леммы~\ref{lem:P1}) сертификат
возвращает $1{,}619026$ --- расхождение равно техническому зазору $10^{-6}$.
Статус метода в размерностях $8$--$9$ --- \emph{численный}: вершины
перечисляются в плавающей точке. Каждая проверяемая величина --- значение
явной рациональной квадратики в решении явной системы линейных уравнений,
однако полное точное перечисление вершин кусков в размерности $8$ пока не
выполнено (\S\ref{sec:open}); в размерности $6$ оно выполнено, и это даёт
второй, независимый точный путь к $\chi(\R^7)\le1029$
(замечание~\ref{rem:r7second}).

\subsection{Ступени ламинирования в $\R^9$: $9604$ и $7203$}\label{ssec:cand9}

Базой служит эйзенштейнова раскраска $E_8/2401$ с $d=\sqrt{7/6}$; в
нормировке $\lambda_1(E_8)^2=3$ имеем $\diam E_8=\sqrt6$ и
$D\bigl((3+\omega)E_8\bigr)=\sqrt7$ (теорема~\ref{thm:eis}). По
лемме~\ref{lem:P1} $\diam\Lambda\le\sqrt{6+t^2}$, по лемме~\ref{lem:P2}
горизонтальные векторы дают $D\ge\sqrt7$; значит, при $t\le1$ строгий диаметр
не превосходит $\sqrt7$, и всё сводится к слоевым векторам.

\emph{Индекс $9604$.} При $m=4$, высоте $t=0{,}93$, сдвиге, смещённом с орбит
глубоких дыр $E_8$, и склейке $a=(5,1,4,5)$ все слоевые векторы дают
$D\ge\sqrt7$, так что $D_{\min}=\sqrt7$; лемма~\ref{lem:P1} даёт
$\diam\le2{,}6200954$, а кусочный сертификат --- $\diam\le2{,}500410$, откуда
вычисленная ширина $1{,}058126$. Сам индекс $9604$ доказан в~\cite{Bounds}
продуктовым исчислением с шириной $\sqrt{63/61}=1{,}0163$; вклад
ламинирования в этой строке --- именно ширина, и она численная. Проверка:
индекс подтверждён точным целочисленным определителем; перечислением
Финке--Похста получены все $266$ векторов подрешётки нормы $\le2\diam$, из них
$240$ горизонтальных закрыты леммой~\ref{lem:P2}, а $26$ слоевых пересчитаны
двумя независимыми процедурами проекции (Дейкстра и SLSQP), согласие ---
до одиннадцатого знака.

Существенны две детали. Во-первых, у $E_8$ \emph{две} орбиты глубоких дыр, и
выбор орбиты меняет результат (при индексе $12005$ они дают $1{,}0397$ против
$0{,}8863$), поэтому сдвиг надо хранить явно. Во-вторых, главный рычаг --- сам
сдвиг: смещение с орбит глубоких дыр разрушает <<дважды глубокую дыру>> (точку,
глубокую сразу для двух соседних слоёв) и опускает настоящий диаметр ниже
границы леммы~\ref{lem:P1}.

\emph{Индекс $7203$ --- доказанная оценка.} Конфигурация $7203=3\cdot2401$
($m=3$, $t\approx1{,}15$) упиралась ровно в одно число --- диаметр, зажатый
между оценкой снизу $2{,}602570$ и границей $2{,}706012$ из леммы~\ref{lem:P1}
при пороге $\sqrt7=2{,}645751$ строго внутри; априорной границы не хватает,
поскольку $R_{\text{базы}}^2+t^2/4=1{,}8306>7/4$. Кусочный сертификат закрыл
вилку численно (сертифицированный диаметр $2{,}602571007$, ширина
$1{,}016591$), но вершины кусков считались в плавающей точке. Точный
рациональный сертификат получен позже, прямым путём: после unimodular-замены
базиса перечисление вершин девятимерной ячейки доводится до конца
($1\,654\,230$ вершин при $752$ фасетах), и в дробях
$\diam^2V_0=1119552292542693/165294140000000<7$, откуда
$d=1{,}0166127\ldots$ Это теорема~5 статьи~\cite{Bounds}; здесь она нужна лишь
как ступень лестницы. Горизонтальный пол $D_{\min}=\sqrt7$ --- теорема
(предложение~\ref{prop:planar} плюс лемма~\ref{lem:P2}), $240$ горизонтальных
векторов закрыты ею.

\subsection{Лестница ширин в $\R^9$}\label{ssec:stair9}

Изначально лестница строилась подбором высоты, при которой слоевые векторы
\emph{ровно} проходят (тогда связывает горизонтальный потолок и
$d=\sqrt7/\sqrt{6+t^2}$ --- равенство). Кусочный сертификат подтвердил все
прежние ступени напрямую и добавил новые, причём ступень $9604$ с
сертифицированным диаметром обгоняет по ширине ступени $12005$ и $14406$ ---
те становятся доминированными (больше цветов при меньшей ширине) и из лестницы
выпадают. Нижняя ступень доказана точно (её ширина --- в дробях); остальные
строки \emph{численные}: ширины вычислены кусочным сертификатом в плавающей
точке.

\begin{center}\small
\begin{tabular}{r r r l c}
\toprule
$k$ & $m$ & $t$ & $d$ & \\
\midrule
$7203$  & $3$  & $1{,}15$ & $1{,}0166127$ & \stC\\
$9604$  & $4$  & $0{,}93$ & $1{,}058126$ & \stN\\
$16807$ & $7$  & $0{,}45$ & $1{,}062345$ & \stN\\
$19208$ & $8$  & $0{,}39$ & $1{,}066688$ & \stN\\
$24010$ & $10$ & $0{,}33$ & $1{,}070453$ & \stN\\
$28812$ & $12$ & $0{,}27$ & $1{,}073621$ & \stN\\
\midrule
--- & $\infty$ & $\to0$ & $\to\sqrt{7/6}=1{,}080123$ & \stP\\
\bottomrule
\end{tabular}
\end{center}

\noindent Предел лестницы содержателен: при $m\to\infty$ высота $t\to0$,
решётка вырождается в $E_8\times\R$, и ширина стремится к $\sqrt{7/6}$ ---
ровно к ширине самой конструкции $E_8/2401$ (теорема~\ref{thm:eis}).
Ламинирование непрерывно интерполирует между рекордом $\R^8$ и $\R^9$.

\subsection{Башня ламинаций $E_8$ при $n\ge10$}\label{ssec:tower}

Ламинирование $E_8$ с сохранением $\lambda_1$ удерживает $\rho<2$ во всех
размерностях. Высота при этом фиксирована условием $R^2+t^2\ge1$, и рекурсия
леммы~\ref{lem:P1} даёт $\diam^2_{n+1}\le\tfrac34\diam^2_n+1$ --- неподвижная
точка~$4$, а старт $\diam^2_8=2$ ниже неё, так что $\rho<2$ и $d>1$ при всех
$n\ge8$ для раскраски $\Gamma=3\Lambda$ (предложение~\ref{prop:mult}). Сам
индекс $3^n$ --- классический ориентир Лармана--Роджерса; новизна здесь в
явных \emph{ширинах} интервалов и в единообразии конструкции. Вычисленные
ширины башни: $1{,}1795$ при $n=10$, $1{,}1166$ при $n=11$ и $1{,}0755$ при
$n=12$ (статус численный: неравенством в цепочке является только рекурсия для
радиуса покрытия; $\lambda_1=1$ проверяется напрямую, а $D_{\min}=2$ --- по
предложению~\ref{prop:mult} и контрольному перечислению всех векторов
$3\Lambda$ нормы $\le2\diam$). В размерности $12$ решётка Коксетера--Тодда даёт
доказанную ширину $\sqrt{3/2}=1{,}2247$ при том же индексе (\S\ref{ssec:k12});
измеренные $\rho$ башни при $n=10,11$ ($\sqrt{8/3}$ и $\sqrt3$) подсказывают,
что и там строгие ширины должны подняться до $1{,}2247$ и $1{,}1547$ --- для
этого нужен сертифицированный радиус покрытия $\Lambda_{10}$, $\Lambda_{11}$.

Отметим ловушку, на которой легко получить внешне строгую, но неверную
рекурсию: у радиуса покрытия здесь \emph{две} роли, и подставлять в них одну и ту
же оценку нельзя. Высоту $t^2=1-R^2$ надо брать из \emph{нижней} границы $R$
(иначе слоевой вектор оказывается короче единицы и $\lambda_1$ незаметно падает),
а диаметр распространять по \emph{верхней}. Замкнутая форма
$\diam^2_n=4-2(3/4)^{n-8}$ строга лишь до $n=10$, где радиус покрытия
$\Lambda_9$ известен точно.

\subsection{Слоевое ламинирование в $\R^{10}$: кандидат индекса 28812}
\label{ssec:dim10lam}

За пределом границы $\diam\le2\sqrt{R_0^2+R_L^2}$ измеренный диаметр
оказывается заметно меньше её, и индекс падает: берутся те же два блока, что в
оценке $45619$ работы~\cite{Bounds}, но слои $A_2$ сдвигаются внутрь ячейки
$E_8$ $\Z[\omega]$-эквивариантными поправками, и слоевое разделение занимает у
поправок, а не у масштаба слоя. Как было и с $7203$, всё упирается в одно
число --- верхнюю границу диаметра, и её даёт обобщение кусочного сертификата
на двумерную слоевую решётку. Путь, которым в итоге закрыли $7203$ --- прямое
перечисление вершин ячейки, --- здесь не годится: там их полтора миллиона уже
в размерности~$9$.

Прямое обобщение потребовало бы общего измельчения \emph{трёх} сдвинутых
разбиений Вороного базы --- в размерности~$8$ это перечисление вершин по
${\sim}10^5$ кускам. Этого удаётся избежать редукцией. Для точки $z$ в
треугольнике Делоне слоевой решётки (вершины $v_k$, описанный радиус $R_L$,
сторона $s$) минимум по трём слоям не больше минимума по двум ближайшим к $z$
вершинам; для пары $(k,l)$ задача расщепляется: вдоль ребра $v_kv_l$ это ровно
одномерная двухслойная задача \S\ref{ssec:cert} с высотой $s$ и сдвигом
$c_l-c_k$, а поперечное отклонение на области, где именно эти две вершины
ближайшие (треугольник $v_kv_l\,o$ с центром описанной окружности $o$), не
превосходит инрадиуса $R_L/2$. Отсюда
\[
R(\Lambda)^2\ \le\ \max_{(k,l)}\ \mathrm{TwoLayer}(c_l-c_k,\,s)\ +\ R_L^2/4 .
\]
Редукция ничего не теряет: при нулевых сдвигах она возвращает в точности
продуктовое значение $R_0^2+R_L^2$. По $\Z[\omega]$-эквивариантности три
разности сдвигов треугольника --- единичные кратные одной
($g$, $\omega g$, $-\omega^2 g$), то есть изометричные копии, и достаточно
одного вызова двухслойного сертификата; дополнительно интервал высот режется
на полосы, в каждой из которых допустимое окно $z$ сужается по мере удаления
от ребра.

Для конструкции индекса $28812=2401\cdot12$ (масштаб слоя $a=0{,}95$)
сертифицированный диаметр равен $2{,}535952$ при $D_{\min}=\sqrt7$, откуда
вычисленная ширина равна $1{,}043297$: по числу цветов это в $2{,}05$ раза
меньше $3^{10}$ и на $37\%$ меньше доказанного $45619$, причём и ширина
$1{,}0433$ больше ширины $\sqrt{217/214}=1{,}0070$ произведения. Горизонтальный
пол $D_{\min}=\sqrt7$ --- теорема (предложение~\ref{prop:planar} плюс аналог
леммы~\ref{lem:P2}: $V_0(\Lambda)\subseteq V_0(E_8)\times\R^2$); слоевые
векторы и вершины кусков сертификата считаются в плавающей точке, поэтому
конструкция остаётся \emph{численным кандидатом} --- в том же положении, в
каком до появления точного сертификата был $7203$. Статусы соседних
индексов по тому же сертификату:

\begin{center}\small
\begin{tabular}{r r r l}
\toprule
индекс & $N$ & $a$ & сертификат полос (плавающая точка)\\
\midrule
$28812$ & $12$ & $0{,}95$ & $d\ge1{,}043297$ --- проходит\\
$28812$ & $12$ & $1{,}10$ & $d\ge1{,}002867$ --- проходит\\
$21609$ & $9$  & $1{,}15$ & $d\ge1{,}000477$ --- см.\ ниже\\
$21609$ & $9$  & $1{,}05$ & $d\ge0{,}990348$ --- не проходит\\
\bottomrule
\end{tabular}
\end{center}

Индекс $21609=3^2\cdot7^4$ (в $2{,}73$ раза ниже $3^{10}$) мы сознательно
не считаем даже кандидатом того же уровня: его сертификат формально проходит,
но запас $4{,}8\cdot10^{-4}$ --- величина, при которой мы не готовы опираться
на плавающую точку (хотя он и на три порядка больше калибровочного зазора
$10^{-6}$); рационализация здесь обязательна прежде любого утверждения.
Контроль на артефакт занижения измеренного радиуса у всех конфигураций один и
тот же: при увеличении числа восхождений с $400$ до $5000$ диаметр меняется
меньше чем на $3\cdot10^{-4}$.

\subsection{Второй путь к $\chi(\R^7)\le1029$}\label{ssec:r7second}

Точный сертификат оценки $\chi(\R^7)\le1029$ в~\cite{Bounds} перечисляет все
$30\,368$ вершин семимерной ячейки Вороного. Кусочный сертификат даёт второй
путь, не разделяющий с первым ни одной вычислительной детали, --- и в
размерности базы $6$ его удаётся провести целиком в точной арифметике.

\begin{remark}[независимая точная проверка $1029$]\label{rem:r7second}
В мажоранту $\varphi=A_0+(A_1-A_0+t^2)^2/(4t^2)$ высота слоя входит только как
$t^2$: иррациональность исчезает, и на каждом куске
$P_s=V_0^{L}\cap(s+V_0^{L})$ общего измельчения двух сдвинутых разбиений
Вороного \emph{базы} функция $\varphi$ остаётся выпуклой квадратикой, максимум
которой берётся в вершинах, а вершины считаются точным перечислением. Полнота
списка кусков доказуема в строку: $|s|\le|x|+|x-s|\le2R_{\mathrm{cov}}(L)$.
Существенно, что любое \emph{ослабление} безопасно --- подмножество неравенств
куска даёт объемлющий многогранник, то есть оценку сверху, --- поэтому
плавающая точка используется лишь как оракул выбора очередного неравенства и на
корректность не влияет.

Для базы $E_6^*$ ($126$ опорных полупространств) получаются $167$ кусков, из них
$103$ доказано пустыми, и
\[
\max_s\max_{P_s}\varphi=\frac{103310189182571717}{59047277850000000}
=1{,}7496\ldots<\frac74,
\]
откуда $\diam^2V_0\le6{,}99847\ldots<7=D_{\min}^2$ и снова
$\chi(\R^7,[1,1])\le1029$. Контроль согласованности сходится: мажоранта не ниже
известного точного $R^2=1{,}640362\ldots$, как и обязана. Сертификат
из~\cite{Bounds} работает с семимерной ячейкой; этот работает в
\emph{шестимерной} базе и семимерную ячейку не строит вовсе
(\path{audit-data/chromatic_research/campaigns/dim7_1029_layer_cert.py}).
\end{remark}
